

%%%%%%%%%%%%%%%%%%%%%%%%%%%%%%%%%%%%%%%%%%%%%%%%%%%%%%%%%%%%%%%%%%%%%%%%
%
% Specify "DGaO-Proc" as document class for your paper.
% Supported options: notcfonts, german (see text below)
%
%%%%%%%%%%%%%%%%%%%%%%%%%%%%%%%%%%%%%%%%%%%%%%%%%%%%%%%%%%%%%%%%%%%%%%%%

\documentclass{DGaO-Proc}

%%%%%%%%%%%%%%%%%%%%%%%%%%%%%%%%%%%%%%%%%%%%%%%%%%%%%%%%%%%%%%%%%%%%%%%%
%
% Temporary solution for ifpdf-redefinition error and wrong paper size
%
%%%%%%%%%%%%%%%%%%%%%%%%%%%%%%%%%%%%%%%%%%%%%%%%%%%%%%%%%%%%%%%%%%%%%%%%

\let\ifpdf\relax
\setlength{\paperwidth}{210mm}
\setlength{\paperheight}{297mm}

%%%%%%%%%%%%%%%%%%%%%%%%%%%%%%%%%%%%%%%%%%%%%%%%%%%%%%%%%%%%%%%%%%%%%%%%
%
% Add any additional packages here.
%
%%%%%%%%%%%%%%%%%%%%%%%%%%%%%%%%%%%%%%%%%%%%%%%%%%%%%%%%%%%%%%%%%%%%%%%%

%\usepackage{}

%%%%%%%%%%%%%%%%%%%%%%%%%%%%%%%%%%%%%%%%%%%%%%%%%%%%%%%%%%%%%%%%%%%%%%%%
%
% Specify a year. This will be printed in the footer on each page.
%
%%%%%%%%%%%%%%%%%%%%%%%%%%%%%%%%%%%%%%%%%%%%%%%%%%%%%%%%%%%%%%%%%%%%%%%%

\year{{{YEAR}}}

%%%%%%%%%%%%%%%%%%%%%%%%%%%%%%%%%%%%%%%%%%%%%%%%%%%%%%%%%%%%%%%%%%%%%%%%
%
% Specify your title.
%
%%%%%%%%%%%%%%%%%%%%%%%%%%%%%%%%%%%%%%%%%%%%%%%%%%%%%%%%%%%%%%%%%%%%%%%%

\title{Document Template for DGaO Proceedings}

%%%%%%%%%%%%%%%%%%%%%%%%%%%%%%%%%%%%%%%%%%%%%%%%%%%%%%%%%%%%%%%%%%%%%%%%
%
% Specify the author(s). Use asterisks ('*') like in the booklet to
% denote different affiliations.
%
%%%%%%%%%%%%%%%%%%%%%%%%%%%%%%%%%%%%%%%%%%%%%%%%%%%%%%%%%%%%%%%%%%%%%%%%

\author{Peter Primary*, Susan Secondary{**}}

%%%%%%%%%%%%%%%%%%%%%%%%%%%%%%%%%%%%%%%%%%%%%%%%%%%%%%%%%%%%%%%%%%%%%%%%
%
% Specify the affiliations of the authors. Use 'newline' (\\) after
% each entry.
%
%%%%%%%%%%%%%%%%%%%%%%%%%%%%%%%%%%%%%%%%%%%%%%%%%%%%%%%%%%%%%%%%%%%%%%%%

\affiliation{*Institute for This and That, University of Somewhere\\{**}Secondary Ltd., Anytown}
	
%%%%%%%%%%%%%%%%%%%%%%%%%%%%%%%%%%%%%%%%%%%%%%%%%%%%%%%%%%%%%%%%%%%%%%%%
%
% Specify the contact mail address of the main author.
% (automatically prefixed with "mailto:")
%
%%%%%%%%%%%%%%%%%%%%%%%%%%%%%%%%%%%%%%%%%%%%%%%%%%%%%%%%%%%%%%%%%%%%%%%%
	
\contactmail{peter.primary@somewhere.edu}

%%%%%%%%%%%%%%%%%%%%%%%%%%%%%%%%%%%%%%%%%%%%%%%%%%%%%%%%%%%%%%%%%%%%%%%%
%
% Specify a short abstract.
%
%%%%%%%%%%%%%%%%%%%%%%%%%%%%%%%%%%%%%%%%%%%%%%%%%%%%%%%%%%%%%%%%%%%%%%%%

\abstract{This is a short abstract which is at most five lines long and not 
identical to the longer abstract that has been submitted during registration. 
The following text moves down if the abstract gets longer.}

\bibliographystyle{osajnl}

%%%%%%%%%%%%%%%%%%%%%%%%%%%%%%%%%%%%%%%%%%%%%%%%%%%%%%%%%%%%%%%%%%%%%%%%
%%%%%%%%%%%%%%%%%%%%%%%%%%%%%%%%%%%%%%%%%%%%%%%%%%%%%%%%%%%%%%%%%%%%%%%%
%%%%%%%%%%%%%%%%%%%%%%%%%%%%%%%%%%%%%%%%%%%%%%%%%%%%%%%%%%%%%%%%%%%%%%%%

% Here your document starts !!

\begin{document}

\section{Introduction}

This document class for \LaTeX\ contains the standard format for the 
proceedings of the \dgao\ conference. To ensure a common appearance for all articles, 
we ask you to adopt the formatting as it is defined here.

The entire article may not be longer than two pages. Longer articles get rejected.

\section{Fonts}

The text is typeset in ``Helvetica 10pt'' for optimal readability on the screen. 
As default, {\tt DGaO-Proc} automatically includes the {\tt tc}-Fonts (Text Companion), 
which provide additional symbols like \um. If you don't have the
{\tt tc}-Fonts installed on your system, use the {\tt notcfonts} option. If you decide
to use the {\tt tc} fonts, please do not use an ugly construct like $\mu$m.

The symbol \permil\ is not taken from the {\tt tc} fonts, as there is no equivalent glyph
for sans serif fonts. You may use the symbol \na\ for ``numerical aperture''. These symbols
work both in text and math mode.

\section{Header}

The template provides commands to typeset the header: \verb|\year|, \verb|\title|, \verb|\author|,
\verb|\affiliation|, \verb|\contactmail|, and \verb|\abstract|. Please see the \TeX\ source of this 
document on how to use them. 

If the authors belong to different organizations 
the affiliation is marked with an asterisk (*). Each article should have only one contact mail address.
The tag {\tt mailto:} is inserted automatically to make it a clickable link in the PDF document.

\section{Figures}

Figures are included using the
\verb|\columnfigure| comand. It is used like this:

\verb|\columnfigure[width]{file}{caption}|

Figures do not float. They are included in the text at the same place where you put them.

\columnfigure[60mm]{DGAOLOGO}{The \dgao\ logo. Images should not be wider than one column. (The width is an optional parameter.)}

The parameter {\tt width} specifies the horizontal scaling of the image. 
If possible please scale figures so that they are not wider than one column. 
This command also automatically sets a label which consists of {\tt fig:} followed by the filename. 
For each figure there should be a reference in the text. 
See Fig.~\ref{fig:DGAOLOGO} for an example.

Please use vector graphic formats (like EPS) for diagrams or sketches. Of course you can use
bitmaps (e.\,g.\ PNG), but this may affect the quality of the image. The images may be in full 24\,bit color. 

\section{Tables}

Tables are typeset using the {\tt columntable} environment. It is used like this:

\verb|\begin{columntable}{spec}{caption}{label}|\\
\verb|\end{columntable}|

% Of course you may define your own commands inside the TeX source
\newcommand{\st}{\superscript{\scriptsize st}}
\newcommand{\nd}{\superscript{\scriptsize nd}}
\newcommand{\rd}{\superscript{\scriptsize rd}}

\begin{columntable}{|l|l|l|l|}{Comparison of results}{comparison}
\hline
Method & 1\st\ Result & 2\nd\ Result & 3\rd\ Result\\
\hline
Old & \dots & \dots & \dots\\
\hline
New & \dots & \dots & \dots\\
\hline
\end{columntable}

Here, {\tt spec} is the column specification (like the {\tt tabular}-environment),
{\tt caption} is the table caption printed below the table and {\tt label} is
a user-defined label, which you can use for cross-references.
See Tab.~\ref{tab:comparison} for an example.

\section{Bibliography}

You may use standard \LaTeX\ bibliographies, by applying the {\tt thebibliography} environment.
The entries of the bibliography are numbered in brackets and cited like this~\cite{Companion04}.

\subsection{Citing \dgao\ Proceedings}

Citations of articles published with the \dgao\ Proceedings need to include the \emph{year} of
publication and the \emph{number} of the presentation (e.\,g.\ A13, P45). Do \emph{not} include the 
page number of the abstract in the conference booklet. The presentation number serves as ``page number''.
You may include the URL \verb|http://www.dgao-proceedings.de|. Do \emph{not} link the URL of a PDF
file directly, since the file location may change in the future. Examples (see also ~\cite{Richter05}):

Short version:\\
{\footnotesize
[1] C.~Richter and G.~H\"ausler, ``Highly improved measurement speed of white light interferometry,''
Proc.\ \dgao, p.\ A4 (2005).
}

Long version:\\
{\footnotesize
[1] C.~Richter and G.~H\"ausler, ``Highly improved measurement speed of white light interferometry,''
106\superscript{th} conference of the \dgao, Wroc{\l}aw (Poland), p.\ A4 (2005). 
URL \texttt{http://www.dgao-proceedings.de}
}

The URL is turned into a clickable link in the PDF document.

\subsection{Using Bib\TeX}

You may also use Bib\TeX. The \dgao\ Proceedings use the OSA style (\texttt{osajnl.bst}) for bibliographies.

\section{Web links}

If you want to insert web links into your text, please mark them accordingly, e.\,g.\ with a
{\tt typewriter} font. Always write the complete web address including the procotol tag,
e.\,g.\ \verb|http://www.dgao.de|. This is necessary because web links in the text are
automatically converted to "`clickable"' links later.

\section{Submit your article}

Please submit your article by email to:

\textit{mailto:submit@dgao-proceedings.de}

This email should have the following content:

\begin{enumerate}
\item\textit{Number of your talk or poster}, as it is displayed in the printed catalogue 
(alternatively, you can include the title of the article and the name of the main author).
\item In the email attachment: your article (preferably packed as .zip file).
\end{enumerate}

You may submit your article as PostScript or PDF document. 

\section{Copyright Agreement}

Please fill out the ``Copyright Agreement'' form, sign it, and send the signed form via e-mail to:\\
submit@dgao.proceedings.de

Deutsche Gesellschaft f\"ur angewandte Optik e.\,V.\\
Technische Universit\"at Ilmenau\\
Fachgebiet Technische Optik\\
Postfach 100565\\
98684 Ilmenau\\
Germany

\textbf{\underline{Important:} By signing this form, you agree that the DGaO may publish 
your article. Until we do not have this agreement, your article cannot be published!}

\section{Support}

If you have questions, feel free to contact:

\texttt{mailto:sekretariat@dgao.de}

%%%%%%%%%%%%%%%%%%%%%%%%%%%%%%%%%%%%%%%%%%%%%%%%%%%%%%%%%%%%%%%%%%%%%%%%
%
% Your bibliography. You may also use BibTeX instead.
%
%%%%%%%%%%%%%%%%%%%%%%%%%%%%%%%%%%%%%%%%%%%%%%%%%%%%%%%%%%%%%%%%%%%%%%%%

\iffalse
\begin{thebibliography}{99}
\small

\bibitem{author}A. Author, ``My Great Paper,'' in \textit{Handbook of the Greatest Papers of the 
World}, H.~Superman ed. (Great Publisher, 1909), pp. 156--187.

\bibitem{somoneelse}S. Someone and E. Else, \textit{Our nicest Results in Color} 
(Oh'Really, Cambridge 2001).

\bibitem{whowasthis}W. Whowasthis, ``I do not know anything,'' \textit{Jour.\ Nonsense} 
\textbf{45}(3), 123--13 (2000).

\end{thebibliography}
\fi

\small
\bibliography{mybib}

%%%%%%%%%%%%%%%%%%%%%%%%%%%%%%%%%%%%%%%%%%%%%%%%%%%%%%%%%%%%%%%%%%%%%%%%
%%%%%%%%%%%%%%%%%%%%%%%%%%%%%%%%%%%%%%%%%%%%%%%%%%%%%%%%%%%%%%%%%%%%%%%%
%%%%%%%%%%%%%%%%%%%%%%%%%%%%%%%%%%%%%%%%%%%%%%%%%%%%%%%%%%%%%%%%%%%%%%%%

\end{document}
